\documentclass[12pt,a4paper]{article}

\usepackage[T1]{fontenc}
\usepackage[utf8]{inputenc}
\usepackage[polish]{babel}
\usepackage{lmodern}
\usepackage{microtype}
\usepackage{amsmath,amssymb,mathtools}
\usepackage{geometry}
\usepackage{enumitem}
\usepackage{hyperref}

\geometry{margin=2.5cm}
\setlength{\parindent}{0pt}
\setlength{\parskip}{0.5em}
\numberwithin{equation}{section}
\setlist[enumerate]{itemsep=0.25em,topsep=0.3em}
\setlist[itemize]{itemsep=0.2em,topsep=0.3em}

\newcommand{\R}{\mathbb{R}}
\newcommand{\N}{\mathbb{N}}
\newcommand{\dd}{\,\mathrm{d}}
\newcommand{\zadanie}[1]{\subsection*{Zadanie #1}}
\newcommand{\tresc}{\paragraph{Treść zadania.}}
\newcommand{\teoria}{\paragraph{Potrzebna teoria.}}
\newcommand{\rozw}{\paragraph{Rozwiązanie.}}

\title{Analiza matematyczna i astronomia\\Lista 4 -- rozwiązania}
\author{}
\date{}

\begin{document}
\maketitle

\section*{Zadania z listy L4}

\zadanie{1}
\tresc
Wykazać, że dla każdego \(c\in\R_+\) szereg o wyrazie ogólnym
\[
  a_n=\frac{c^n}{n!},\qquad n\geq 1,
\]
jest zbieżny. Wywnioskować, że \(\frac{c^n}{n!}\to 0\).

\teoria
Kryterium d'Alemberta: jeżeli \(a_n>0\) i
\[
  \lim_{n\to\infty}\frac{a_{n+1}}{a_n}=L<1,
\]
to szereg \(\sum_{n=1}^{\infty}a_n\) jest zbieżny. Dla każdego
szeregu zbieżnego jego wyrazy dążą do zera.

\rozw
Dla \(c>0\) obliczamy iloraz kolejnych wyrazów:
\[
  \frac{a_{n+1}}{a_n}
  =\frac{c^{n+1}}{(n+1)!}\frac{n!}{c^n}
  =\frac{c}{n+1}
  \xrightarrow[n\to\infty]{}0.
\]
Granica jest mniejsza od \(1\), więc z kryterium d'Alemberta szereg
jest zbieżny. Gdy \(c=0\), wszystkie wyrazy dla \(n\geq1\) są równe
zeru, a wniosek pozostaje prawdziwy. Z koniecznego warunku zbieżności
szeregu otrzymujemy
\[
  \boxed{\lim_{n\to\infty}\frac{c^n}{n!}=0.}
\]

\zadanie{2}
\tresc
Niech \(k\in\N\) i \(c>1\). Wykazać, że szereg o wyrazie ogólnym
\[
  a_n=\frac{n^k}{c^n},\qquad n\geq1,
\]
jest zbieżny. Wywnioskować, że \(\frac{n^k}{c^n}\to0\).

\teoria
Zastosujemy kryterium d'Alemberta: wystarczy, by granica ilorazu
\(a_{n+1}/a_n\) była mniejsza od \(1\). Wyrazy zbieżnego szeregu
muszą dążyć do zera.

\rozw
Dla \(n\geq1\) mamy
\[
  \frac{a_{n+1}}{a_n}
  =\frac{(n+1)^k}{c^{n+1}}\frac{c^n}{n^k}
  =\frac{1}{c}\left(1+\frac{1}{n}\right)^k
  \xrightarrow[n\to\infty]{}\frac{1}{c}<1.
\]
Zatem \(\sum_{n=1}^{\infty} n^k/c^n\) jest zbieżny, a w szczególności
\[
  \boxed{\lim_{n\to\infty}\frac{n^k}{c^n}=0.}
\]

\zadanie{3}
\tresc
Wykazać, że szereg
\[
  \sum_{n=1}^{\infty}\frac{1}{4n^2-1}
\]
jest zbieżny i obliczyć jego sumę.

\teoria
Rozkład na ułamki proste może ujawnić teleskopową strukturę szeregu.
Wtedy sumę oblicza się najpierw dla sum częściowych, a następnie
przechodzi z \(N\) do nieskończoności.

\rozw
Rozkładamy mianownik:
\[
  4n^2-1=(2n-1)(2n+1),
  \qquad
  \frac{1}{4n^2-1}
  =\frac12\left(\frac{1}{2n-1}-\frac{1}{2n+1}\right).
\]
Suma częściowa jest więc teleskopowa:
\begin{align*}
  S_N
  &=\sum_{n=1}^{N}\frac{1}{4n^2-1}\\
  &=\frac12\left[
      \left(1-\frac13\right)
      +\left(\frac13-\frac15\right)
      +\cdots
      +\left(\frac{1}{2N-1}-\frac{1}{2N+1}\right)
    \right]\\
  &=\frac12\left(1-\frac{1}{2N+1}\right).
\end{align*}
Stąd
\[
  \boxed{\sum_{n=1}^{\infty}\frac{1}{4n^2-1}
    =\lim_{N\to\infty}S_N=\frac12.}
\]
Szereg jest zbieżny, ponieważ jego sumy częściowe mają skończoną
granicę.

\zadanie{4}
\tresc
Dla \(c>0\) zbadać zbieżność szeregu
\[
  \sum_{n=1}^{\infty}a_n,
  \qquad
  a_n=\frac{c^n n!}{n^n}.
\]
Rozpatrzyć osobno przypadki \(c<e\), \(c=e\) oraz \(c>e\).

\teoria
Kryterium d'Alemberta rozstrzyga przypadki, w których granica
\(a_{n+1}/a_n\) jest różna od \(1\). W przypadku granicy równej \(1\)
potrzebne jest dodatkowe oszacowanie; użyjemy wzoru Stirlinga
\[
  n!\sim \sqrt{2\pi n}\left(\frac{n}{e}\right)^n.
\]

\rozw
Wszystkie wyrazy są dodatnie. Iloraz kolejnych wyrazów wynosi
\[
  \frac{a_{n+1}}{a_n}
  =\frac{c^{n+1}(n+1)!}{(n+1)^{n+1}}
    \frac{n^n}{c^n n!}
  =c\left(\frac{n}{n+1}\right)^n
  \xrightarrow[n\to\infty]{}\frac{c}{e},
\]
bo \(\left(1+\frac1n\right)^n\to e\).
\begin{itemize}
  \item Jeśli \(0<c<e\), to \(c/e<1\). Kryterium d'Alemberta daje
    zbieżność szeregu.
  \item Jeśli \(c>e\), to \(c/e>1\). Istnieje \(q>1\), dla którego
    od pewnego miejsca \(a_{n+1}/a_n>q\). Wyrazy nie dążą do zera
    (w istocie rosną co najmniej geometrycznie), więc szereg jest
    rozbieżny.
  \item Jeśli \(c=e\), to kryterium d'Alemberta nie rozstrzyga.
    Ze wzoru Stirlinga:
    \[
      a_n=\frac{e^n n!}{n^n}
      \sim \sqrt{2\pi n}.
    \]
    W szczególności \(a_n\not\to0\), zatem szereg jest rozbieżny.
\end{itemize}
Ostatecznie
\[
  \boxed{\text{szereg jest zbieżny wtedy i tylko wtedy, gdy }0<c<e.}
\]

\zadanie{5}
\tresc
Zbadać zbieżność szeregów o podanych wyrazach ogólnych:
\[
  \text{a)}\quad a_n=\ln\left(1+\frac1n\right),\qquad
  \text{b)}\quad a_n=\frac{(\arctan n)^n}{2^n},\qquad
  \text{c)}\quad a_n=\sqrt{n^2+1}-n.
\]

\teoria
\begin{itemize}
  \item Szereg teleskopowy rozpoznajemy, zapisując wyraz jako różnicę
    kolejnych wyrazów pewnego ciągu.
  \item Kryterium Cauchy'ego (pierwiastkowe) mówi, że jeśli
    \(\limsup\sqrt[n]{|a_n|}<1\), to szereg \(\sum a_n\) jest
    zbieżny bezwzględnie.
  \item Zasada porównawcza: szereg o wyrazach nieujemnych jest
    rozbieżny, jeśli jego wyrazy są od pewnego miejsca nie mniejsze
    od wyrazów rozbieżnego szeregu.
\end{itemize}

\rozw
\textbf{a) Szereg \(\displaystyle\sum_{n=1}^{\infty}
\ln\left(1+\frac1n\right)\).}
Ponieważ
\[
  \ln\left(1+\frac1n\right)
  =\ln\left(\frac{n+1}{n}\right)=\ln(n+1)-\ln n,
\]
sumy częściowe teleskopują:
\[
  \sum_{n=1}^{N}\ln\left(1+\frac1n\right)
  =\ln(N+1)\xrightarrow[N\to\infty]{}+\infty.
\]
Zatem szereg jest \(\boxed{\text{rozbieżny}}\).

\medskip
\textbf{b) Szereg \(\displaystyle\sum_{n=1}^{\infty}
\frac{(\arctan n)^n}{2^n}\).}
Wyrazy są dodatnie. Stosujemy kryterium pierwiastkowe:
\[
  \sqrt[n]{a_n}=\frac{\arctan n}{2}
  \xrightarrow[n\to\infty]{}\frac{\pi}{4}<1.
\]
Szereg jest więc \(\boxed{\text{zbieżny bezwzględnie}}\).

\medskip
\textbf{c) Szereg \(\displaystyle\sum_{n=1}^{\infty}
\left(\sqrt{n^2+1}-n\right)\).}
Usuwamy niewymierność z licznika:
\[
  \sqrt{n^2+1}-n
  =\frac{(\sqrt{n^2+1}-n)(\sqrt{n^2+1}+n)}
          {\sqrt{n^2+1}+n}
  =\frac{1}{\sqrt{n^2+1}+n}.
\]
Dla \(n\geq1\), \(\sqrt{n^2+1}\leq\sqrt{2}\,n\), zatem
\[
  \sqrt{n^2+1}-n
  \geq \frac{1}{(1+\sqrt{2})n}
  >\frac{1}{3n}.
\]
Szereg harmoniczny \(\sum 1/n\) jest rozbieżny, więc z kryterium
porównawczego rozbieżny jest również rozważany szereg:
\(\boxed{\text{szereg jest rozbieżny}}\).

\zadanie{6}
\tresc
Zbadać zbieżność szeregu naprzemiennego o wyrazach
\[
  a_n=(-1)^{n-1}\left(\sqrt[n]{3}-1\right),
  \qquad n\geq1.
\]
Czy szereg jest zbieżny bezwzględnie? Wskazówka: porównać
\(\sqrt[n]{3}-1\) z wyrazami szeregu harmonicznego.

\teoria
Kryterium Leibniza: szereg naprzemienny
\(\sum(-1)^{n-1}b_n\) jest zbieżny, jeśli \(b_n\geq0\), \(b_n\to0\)
i ciąg \(b_n\) jest nierosnący. Zbieżność bezwzględną
sprawdzamy przez szereg \(\sum |a_n|\).

\rozw
Oznaczmy
\[
  b_n=\sqrt[n]{3}-1
      =\exp\left(\frac{\ln 3}{n}\right)-1.
\]
Mamy \(b_n>0\), a funkcja
\[
  b(x)=\exp\left(\frac{\ln 3}{x}\right)-1,\qquad x>0,
\]
jest malejąca, ponieważ
\[
  b'(x)=-\frac{\ln 3}{x^2}
         \exp\left(\frac{\ln 3}{x}\right)<0.
\]
Ponadto \(b_n\to0\). Kryterium Leibniza daje zbieżność szeregu
\(\sum a_n\).

Sprawdzamy zbieżność bezwzględną. Dla \(x>0\) zachodzi
\(e^x-1>x\), więc
\[
  |a_n|=b_n
  =e^{(\ln3)/n}-1
  >\frac{\ln3}{n}
  >\frac1n,
\]
bo \(\ln3>1\). Z porównania z szeregiem harmonicznym wynika, że
\(\sum |a_n|\) jest rozbieżny. Zatem
\[
  \boxed{\text{szereg jest zbieżny warunkowo, ale nie bezwzględnie}.}
\]

\zadanie{7}
\tresc
Piłkę upuszczono z wysokości \(h=1\,\mathrm{m}\). Po jakim czasie
znajdzie się w spoczynku, jeżeli przy każdym zderzeniu z ziemią traci
połowę energii kinetycznej? Pominąć opór powietrza i rozmiary piłki.

\teoria
W polu grawitacyjnym energia kinetyczna przy uderzeniu wynosi
\(E_k=\frac12mv^2\). Jeśli po zderzeniu pozostaje połowa energii, to
prędkość po odbiciu jest pomnożona przez \(q=1/\sqrt2\), ponieważ
energia jest proporcjonalna do \(v^2\). Czas swobodnego spadania
z wysokości \(h\) wynosi \(\sqrt{2h/g}\). Kolejne czasy lotu tworzą
szereg geometryczny.

\rozw
Czas pierwszego spadania oznaczmy przez
\[
  t_0=\sqrt{\frac{2h}{g}}.
\]
Prędkość tuż przed pierwszym zderzeniem to \(v_0=\sqrt{2gh}\).
Po każdym odbiciu energia, a więc także wysokość maksymalna,
zmniejsza się dwukrotnie. Prędkość po \(n\)-tym odbiciu wynosi
\[
  v_n=\frac{v_0}{(\sqrt2)^n}.
\]
Czas wznoszenia po \(n\)-tym odbiciu jest równy \(v_n/g\), a całkowity
czas kolejnego lotu (wznoszenie i spadanie) to
\[
  \Delta t_n=\frac{2v_n}{g}
  =2t_0\left(\frac1{\sqrt2}\right)^n,\qquad n\geq1.
\]
Łączny czas wszystkich lotów jest zatem
\begin{align*}
  T
  &=t_0+\sum_{n=1}^{\infty}\Delta t_n\\
  &=t_0\left(1+2\sum_{n=1}^{\infty}
                  \left(\frac1{\sqrt2}\right)^n\right)\\
  &=t_0\left(1+\frac{2/\sqrt2}{1-1/\sqrt2}\right)
    =t_0(3+2\sqrt2).
\end{align*}
Dla \(h=1\,\mathrm{m}\) i \(g\approx9.81\,\mathrm{m/s^2}\):
\[
  \boxed{T=(3+2\sqrt2)\sqrt{\frac{2}{9.81}}\,\mathrm{s}
         \approx 2.63\,\mathrm{s}.}
\]
Piłka osiąga spoczynek dokładnie dopiero w granicy nieskończonej
liczby odbić; \(T\) jest skończonym czasem granicznym tego procesu.

\end{document}
